\begin{abstract}
Atmospheric turbulence degrades images through spatially varying blur and geometric distortion, making reconstruction from a single observation ambiguous.
We propose a physics-guided recursive diffusion framework for single-image turbulence restoration that incorporates a differentiable turbulence simulator into the restoration process.
At each stage, a conditional diffusion model estimates a clean image from the current degraded state, while a separate prediction branch estimates the latent variables required by the simulator.
The simulator then renders these estimates at a lower turbulence strength to construct the input to the next stage.
During training, re-rendering each stage's estimates at the original strength provides an observation-consistency loss that anchors the recursive process to the input image.
Physical strength and the internal diffusion timestep are represented separately, allowing the restoration model to account for changing degradation levels while retaining its conditional diffusion formulation.
Training proceeds from direct restoration to learning on simulated intermediate states, followed by joint optimization of the unrolled restoration process with gradients through the simulator.
The resulting framework uses the turbulence forward model to define both the transitions between restoration stages and the observation constraints applied during training.
\end{abstract}
